% Imporditud: tools/import_eksperimendid.py
% Allikas: Eesti füüsikaolümpiaadi eksperimendi ülesannete
%   kogu (Rael Kalda, 2023), ülesanne Ü130, lahendus L130.
\setAuthor{EFO žürii}
\setRound{piirkonnavoor}
\setYear{2005}
\setNumber{G E1}
\setDifficulty{4}
\setTopic{geomeetriline optika}
\setVahendid{punktvalgusallikas, nõguslääts hoidjas, mõõtejoonlaud, valge paberileht}
\setPunktid{10}
\setAllikas{Eksperimendi kogu Ü130, lahendus lk 101}
\setLahendusseis{toores}

\prob{Läätse fookuskaugus}
Kasutades paberilehte ekraanina jälgige, millise varju jätab lääts paberile. Visandage valguse intensiivsuse jaotus paberil sõltuvuses punkti kaugusest optilisest peateljest siis, kui läätse kaugus paberist on umbes 10 cm. Leidke läätse fookuskaugus kasutades antud nähtust.

\hint

\solu
Sõltuvalt läätse kaugusest ekraanist võib olla nähtav läätse hoidja tume vari: juhtum (a), vt joonis (2 p). Kui kaugus on suurem, siis on läätse taguse tumedama piirkonna ja fooni vahel heledam piirkond, kuhu jõuavad nii otse tulevad kiired, kui ka läätsest hajunud kiired: juhtum (b), vt joonis (2 p).

Fookuskauguse võib leida nt siis, kui mõõdame juhtumil (b) heleda piirkonna diameetri D = 2R --- siis, kui läätse kaugus ekraanist on l.

r R

fl

F

Leiame fookuskauguse f sarnastest kolmnurkadest

f +l f ld (2 p). = $\Rightarrow$ f= Dd D$-$d

Mõõtmised: Mõõtmised on dokumenteeritud ja nende põhjal on leitud f (1 p); tulemus on tõepärane (1 p). Mõõtmisi on korratud mitu korda (1 p). Hinnatud on mõõteviga (1 p).
\probend
